\scanpage{8}\mspage{6}
\uline{Th. 1.8.}

$f\colon X\longrightarrow Y$ loc. de prés. finie

\margin{\uline{NB} Si on suppose $Y$ loc. noeth, alors l'énoncé équivaut à dire que, si $X$ est réduit, alors au voisinage de $x$ $f$ est normal.}
\margin{Si $X$ et $Y$ sont loc. noeth, on doit pouvoir supprimer l'hyp. que $f$ est \deleted{\unsure{loc.} \unsure{de t.}} de prés finie (moyennant des hypothèses d'excellence éventuellement --}

\deleted{\unsure{\textcircled{1}}} \unsure{Soient|Soit} $x\in X$, $y=f(x)$

a) $Y$ réduit en $y$

b) $f$ univ. ouvert en les gén. maximales de $x$ dans $X_y$

c) Pour tte gén. \deleted{\unsure{maximale}} $y'$ de $y$, et toute \inserted{les} générisation $x'$ de
$x$ appartenant à $X_{y'}$, les $\dim_{x'}X_{y'}$ sont égaux à \deleted{\unsure{$\dim X_y$}} $\dim_x\deleted{\unsure{X}}X_y$

d) $X_y$ \deleted{\unsure{est|réduit} \unsure{$x$} séparable en} réduit en les gén. maximales de $x$ dans $X_y$

e) $(X_y)_{\text{réd}}$ est géom\textsuperscript{t}. normal au voisinage de $x$

% slanted
\margin{peut être remplacé par c') $f$ univ. ouvert dans un voisinage de $x$, et c) \unsure{valable si} $y'$ est maximal}

[d) et e) équivalent à : $\exists$ voisinage ouvert $U$ de $x$ \inserted{dans $X_y$}, tel que
$\overline U$ \inserted{$=U\otimes_{k(y)}\overline{k(y)}$} soit gén. réduit et $\overline U_{\text{réd}}$ normal].

% Conclusion and abandoned formulations, from the centre down.
\uline{Alors} \deleted{\unsure{1°|1}} Il existe un voisinage ouvert $U$ de $x$, tel que
le morphisme \inserted{$U_{\text{réd}}\to Y$ soit \deleted{normal} loc. de présentation finie et}
\deleted{$U\to Y$ soit univ. ouvert, \unsure{et} \unsure{intègre} \unsure{fini} et à fibres géométriques normales, donc \inserted{localement} équidimensionnel}
\deleted{à fibres géométriques \inserted{\unsure{gén.}} réduites telles que $(X_y)_{\text{réd}}$ normal.}

\deleted{2) Si $Y$ est réduit, la condition envisagée est vérifiée,}
\deleted{et si $Y$ est réduit, alors \quad Soit $Y_0=Y_{\text{réd}}$, $X_0=X\times_Y Y_0$.}

\deleted{\uline{Cor.} \textcircled{?} Conditions équivalentes : (i) Les conditions envisagées dans \textcircled{1} sont satisfaites en $\forall\,x\in X$, \unsure{et que $f$ est univ. ouvert} fibre \unsure{géométriquement} \unsure{unib|univ.} \inserted{i.e. (i bis) $f$ est univ. ouvert, \inserted{loc} équidimensionnel, ses fibres sont gén. séparables et ses fibres réduites géom. normales.}}
\deleted{(ii) $\exists$ un sous--}
\deleted{(ii) Il existe un ouvert $V$ de $X_0$, tel que pour $\forall\,y\in Y$, $V_y$ soit dense dans $X_y$, et $V\longrightarrow Y_0$ un morphisme \deleted{normal} séparable. De plus \deleted{$X_0\to Y_0$ est à fibres loc. équidimensionnelles, et} pour $\forall\,y\in Y_0$, $(X_{0y})_{\text{réd}}$ est géom. normal.}
\deleted{(iii) \deleted{Il existe un sous $X_{0\,\text{réd}}\to Y_0$} $X_{\text{réd}}\to Y_{\text{réd}}$ est normal et \inserted{loc.} de présentation finie}
\deleted{(iiii bis) Il existe un sous-schéma \deleted{\unsure{X}} $Z_0$ de $X_0$, \deleted{plat} normal \inserted{et loc. de prés. finie} sur $Y_0$, \deleted{et} ayant \unsure{m\textsuperscript{me}} espace sous-jacent que $X_0$, et coïncidant avec $X_0$ en les pts génériques des fibres \ldots}

normal. De plus, il existe un ouvert \inserted{relativement} schématiquement dense $V$ de
$U$ \deleted{\unsure{du|dans}} $S$, tel que $V\to X$ soit une \unsure{imm.|univ.} ouverte
% archive sheet number: 7
